% Check criterion, deliverable, and the Decision Log question of Day 3.
% The criterion and the deliverable come from the Day 3 section of the
% syllabus and from Labs/day03/README.md.
\begin{frame}{Check criterion}
  \small
  The instructor checks this at the end of the lab:
  \begin{itemize}
    \item \cmark\ The notebook prints \code{complete} for tasks 1, 2, and 4.
    \item \cmark\ The Serial Monitor shows \code{Self-test: PASS}.
    \item \cmark\ The board shows the correct class for each of the four
          motions.
    \item \cmark\ The report has the three predictions: task 3, the arena
          estimate, and the PSRAM question. You wrote them before the
          measurement.
    \item \cmark\ The comparison table has measured numbers for your
          sketch: flash, arena, and latency, with PSRAM and with no PSRAM.
    \item \cmark\ The comparison table has numbers for the path of Edge
          Impulse. Each number says if it is a measurement or an estimate.
    \item \cmark\ The Decision Log gives numbers and names one trade-off.
  \end{itemize}
  \vskip 0.4em
  \textbf{Hand in:} the file \code{report.md}. Show the kit to the
  instructor.
\end{frame}

\begin{frame}{Decision Log}
  Write about 100 words. State one design decision, give your measured
  numbers, and name the trade-off.
  \vskip 0.4em
  \begin{exerciseblock}
    A product must classify the motion of a package for one year with a
    battery. Which path do you select for the firmware?
    \begin{itemize}
      \item Your own sketch with TensorFlow Lite Micro.
      \item The library of Edge Impulse.
    \end{itemize}
  \end{exerciseblock}
  \vskip 0.2em
  A good Decision Log:
  \begin{itemize}
    \item gives the flash, the RAM, and the latency of the two paths from
          your table,
    \item says what the latency means for the time that the processor can
          sleep,
    \item names what you lose with your decision. Example: control of the
          code against the speed of optimized kernels.
  \end{itemize}
\end{frame}
